\section{Worker I: finite-dimensional proofs for all words}
\label{sec:words}

\subsection{Model and certificate}

Let $A$ be a finite alphabet and $d\ge1$. A rational linear word model consists of
an initial row $u\in\Q^{1\times d}$, a matrix $M_a\in\Q^{d\times d}$ for every
$a\in A$, and an output column $v\in\Q^{d\times1}$. For
$w=a_1\cdots a_m$, put
\[
M_w=M_{a_1}\cdots M_{a_m},\qquad M_\epsilon=I,
\qquad f(w)=uM_wv.
\]
The worker decides whether $f(w)=0$ for every finite word. Equality of two models
reduces to this problem by block-diagonal composition:
\[
u=(u_1,u_2),\quad M_a=\operatorname{diag}(M_{1,a},M_{2,a}),\quad
v=\begin{pmatrix}v_1\\-v_2\end{pmatrix}.
\]
This reduction requires the same alphabet and an exact correspondence between
its symbols. It does not identify two source constructors because they have
similar printed names.

A positive certificate contains an arbitrary matrix $B\in\Q^{b\times d}$,
a row $\alpha\in\Q^{1\times b}$, and matrices
$C_a\in\Q^{b\times b}$ satisfying
\begin{equation}\label{eq:word-cert}
u=\alpha B,\qquad BM_a=C_aB\quad(a\in A),\qquad Bv=0.
\end{equation}
The search usually chooses independent rows for $B$, but the checker does
\emph{not} need to prove independence, minimality, or even reachability of those
rows. Any finite simulation satisfying these identities suffices. The rank-zero
case is legitimate: it certifies exactly the zero initial row.

\begin{theorem}[Soundness of the word certificate]\label{thm:word-sound}
If \eqref{eq:word-cert} holds, then $uM_wv=0$ for every $w\in A^*$.
\end{theorem}
\begin{proof}
Induction on the word gives $BM_w=C_wB$, where $C_w$ is the analogous product
of the small matrices. Consequently
$uM_wv=\alpha BM_wv=\alpha C_wBv=0$. The empty word uses the identity matrices.
Every equality is a matrix identity over the same rational field.
\end{proof}

The certificate says more than ``these examples worked''. It provides a finite
state representation that is closed under every transition and whose entire
represented space has output zero. This is the missing universal argument.

\subsection{Search: exact reachable-space closure}

Define the reachable row space
\[
\mathcal R=\Span_\Q\{uM_w:w\in A^*\}.
\]
A breadth-first agenda stores a row basis for this space together with a word
that generates each stored row. The procedure is elementary.

\begin{lstlisting}
check u * v; return the empty word if it is nonzero
insert u if independent; queue its basis index
while the queue is nonempty:
    take one stored row b with generating word w
    for each letter a:
        y := b * M[a]
        if y * v != 0: return the concrete word w ++ [a]
        if y is independent of the stored rows:
            insert y and queue it, retaining w ++ [a]
solve u = alpha * B and B * M[a] = C[a] * B
return B, alpha, C
\end{lstlisting}

The implementation incrementally maintains echelon rows and their exact
coordinates in the original, word-labelled rows. It never uses a floating-point
rank tolerance. A word is checked against the original model, not against the
normalised echelon rows, which need not themselves be reachable by a single
word.

\begin{theorem}[Termination and completeness of the rational fragment]
Without a resource interruption, reachable-space search terminates and returns
either a valid certificate \eqref{eq:word-cert} or a word with nonzero output.
A separating word exists of length at most $d-1$ whenever the universal identity
is false.
\end{theorem}
\begin{proof}
Every admitted row is independent of the previous rows, so at most $d$ rows
are admitted. Each is expanded once under each letter. At agenda exhaustion,
the stored span contains $u$ and is closed under all $M_a$, so it contains every
reachable row. Conversely, every admitted row is reachable. Thus it equals
$\mathcal R$, and the coordinate equations for the certificate are solvable.
Every stored or tested row is checked against $v$; if none separates, the
certificate satisfies $Bv=0$.

For the length bound, let $R_j$ be the span of reachable rows from words of
length at most $j$. If $R_j=R_{j+1}$, it is already transition-closed and all
later spaces agree. Starting with dimension at most one, the ascending chain
must either stabilise earlier or reach its maximum dimension by $j=d-1$.
Therefore $R_{d-1}=\mathcal R$. A nonzero linear functional on $\mathcal R$
is nonzero on at least one of the word generators of $R_{d-1}$.
\end{proof}

This is the standard weighted-automaton argument, specialised to a small
certificate interface~\cite{kiefer}. The new checker deliberately does less
than the search: it multiplies rational matrices and compares the results.
It need not trust how Gaussian elimination was implemented.

\subsection{Complexity, proof size, and useful nonminimal certificates}

Incremental reduction and dense matrix-vector multiplication take $O(d^2)$
field operations per candidate row. There are at most $|A|d$ expansions, so the
basic search is $O(|A|d^3)$ field operations. Dense verification of
\eqref{eq:word-cert} costs
\[
O\bigl(|A|(bd^2+b^2d)+bd\bigr),\qquad b\le d
\]
for certificates emitted by this search. Storage is
$O(bd+|A|b^2)$ rational entries, in addition to the original model.
These are \emph{field-operation} bounds, not constant-time bit bounds.
Intermediate rational numerators and denominators can grow substantially.

In a production search, fraction-free elimination or modular discovery with
exact reconstruction is worth considering. In the acceptance path, however,
reconstruction must be followed by the same exact identities. A modular rank
calculation is not evidence that a rational equality holds. The proposed FLINT
backend is only a faster producer of candidates.

Minimising $B$ is optional. A small nonminimal certificate can be much easier to
find than an absolutely minimal one, and its algebraic checks remain sound.
For example, the Cassini lift below has dimension ten and an emitted reachable
basis of rank eight; no claim of output-minimal automaton construction is made.
A later certificate optimiser can reduce rows if it provides another valid
simulation. Proof size, not minimality terminology, is the practical objective.

\subsection{Infinite input alphabets by polynomial interpolation}
\label{sec:alphabet}

Many relevant folds consume arbitrary integers or rationals, not letters from a
small enumerated datatype. Suppose the transition matrix depends polynomially
on an input $x$:
\begin{equation}\label{eq:poly-transition}
M(x)=\sum_{j=0}^{D}x^jM_j,\qquad M_j\in\Q^{d\times d}.
\end{equation}
A positive certificate checks $BM_j=C_jB$ for every coefficient matrix. Then
\[
BM(x)=\left(\sum_{j=0}^{D}x^jC_j\right)B
\]
for every rational $x$, so Theorem~\ref{thm:word-sound} proves the identity for
every finite sequence of rational inputs. The same positive algebraic argument
transports to a compatible characteristic-zero field through its rational
embedding; the source adapter must prove that transport.

There is also a complete finite search alphabet.

\begin{theorem}[Interpolation alphabet]\label{thm:interp}
For \eqref{eq:poly-transition}, choose $D+1$ distinct rationals
$t_0,\ldots,t_D$. The reachable space generated by all $M(x)$, $x\in\Q$, equals
the reachable space generated by $M(t_0),\ldots,M(t_D)$.
\end{theorem}
\begin{proof}
Every grid transition is an allowed rational transition, giving one inclusion.
The Vandermonde matrix $(t_i^j)_{0\le i,j\le D}$ is invertible, so each
coefficient matrix $M_j$ is a rational linear combination of the grid matrices.
A row space closed under every grid matrix is consequently closed under every
$M_j$ and hence every $M(x)$. This gives the other inclusion.
\end{proof}

The implementation uses the grid $0,1,\ldots,D$. It searches with those actual
transitions, then emits closure actions for the \emph{coefficient} matrices.
A negative result is therefore a concrete list of grid inputs. This detail
matters: directly searching the coefficient matrices could produce a sequence
of coefficient indices that is not a legal sequence of source inputs.

The grid is not an empirical sample. Completeness rests on the proved degree
bound. If that bound is understated, the conclusion is false. The test
\[
s' = s+x(x-1),\qquad s_0=0
\]
is zero for every word over inputs $0$ and $1$, but one input $2$ yields output
$2$. The compiler computes degree two, and the worker returns $[2]$ as a
separating input. This is a regression control against treating a small input
set as universal without checking interpolation hypotheses.

For several input coordinates, a rectangular interpolation grid is a valid
extension, but its size is a product of degree bounds. A sparse multivariate
unisolvent set could improve search. Neither extension is implemented here.

\subsection{Affine state updates and polynomial observations}
\label{sec:lift}

A useful front end need not ask users to write matrices. Let the state be
$z=(z_1,\ldots,z_m)$, let each update be affine in $z$, and let the desired
output relation be a polynomial $H(z)$ of degree at most $e$. Form the feature
row
\[
\Phi_e(z)=(z^\alpha)_{|\alpha|\le e},\qquad
\dim\Phi_e=\binom{m+e}{e}.
\]
Substitution of an affine update into a polynomial of degree at most $e$ does
not increase its state degree. Thus each update $T_a$ has a matrix satisfying
\begin{equation}\label{eq:lift}
\Phi_e(T_a(z))=\Phi_e(z)M_a,
\qquad H(z)=\Phi_e(z)v.
\end{equation}
The initial feature row is $u=\Phi_e(z_0)$. Equations \eqref{eq:lift}, followed
by the matrix certificate, prove $H$ vanishes on every reachable state.

Coefficients of an affine state update may themselves be polynomial in an
input $x$. Substitution then gives matrices of form \eqref{eq:poly-transition}.
The relevant input degree is the degree \emph{after lifting}; it can exceed the
degree in the unlifted update. The supplied compiler expands substitutions,
extracts state monomial coefficients, determines input degree, and checks the
symbolic substitution identities with SymPy. Its positive certificates are
checked in the matrix model. This compiler is not yet an independently proved
Lean reifier.

\begin{example}[Two related Horner folds]\label{ex:horner}
Start at $(h,j,g)=(0,0,0)$ and consume a list through
\[
h'=bh+x,\qquad j'=bj+2x+3,\qquad g'=bg+1,
\]
where $b$ is a fixed rational. The desired identity is $j=2h+3g$ after every
list. The degree-one features $(1,g,j,h)$ give a four-dimensional model with
input degree one. The tested instances $b=-2,0,1,3$ each emit a rank-three
closure certificate after six basis-letter expansions. A direct handwritten
induction is short; the important automatic step is deriving and certifying
the hidden state relation without being given that induction proof.
\end{example}

\begin{example}[A polynomial observer, not a polynomial sequence]
Start with $(a,b,c)=(0,1,1)$ and allow either transition
\[
T_1(a,b,c)=(b,a+b,-c),\qquad
T_2(a,b,c)=(a+b,a+2b,c).
\]
The observer is $H=b^2-ab-a^2-c$. Since $T_2=T_1^2$, words describe arbitrary
blocks of one or two Fibonacci advances, and the identity is the Cassini
relation with the parity sign carried in $c$. Degree-two lifting has ten
features. Exact search emits a rank-eight certificate after sixteen expansions.
This proves an output identity for an infinite collection of transition words;
it does not claim that the Fibonacci values themselves are polynomials in the
word length.
\end{example}

\subsection{What the front end must reject or leave as an obligation}

The first source fragment should be deliberately narrow: fixed rational
coefficients, affine state updates, polynomial observations, and either an
explicit finite transition alphabet or one polynomial input coordinate.
A state update such as $z'=z^2$ is rejected by the supplied compiler. A broader
Koopman-style lift may exist for some nonlinear systems, but a finite closure is
not automatic and must not be presumed from a finite set of experiments.

Natural subtraction, integer division, overflow arithmetic, and tests on state
values do not become rational affine operations merely because their notation
resembles them. Casts require homomorphism proofs; a branch requires an exact
case split or a documented over-approximation. If all branch actions are allowed
in arbitrary order, positive closure is sound for the resulting larger model,
but a negative word may be infeasible in the source program. Such a word is only
guidance until source execution or a proved feasibility argument confirms it.

An unknown or unsupported translation is not a disproof. For a translated
exact model, by contrast, a separating word is genuine positive evidence of an
inequality. The certificate interface distinguishes these cases instead of
routing both through a single ``failed'' result.
